\section{A coherent raw correction under changes of cutoff}
\label{sec:coherent-raw-cutoff}
Increasing the low-frequency region changes both the extracted-edge
convolution and the residual inverse transform. Estimating their
changes separately loses an exact cancellation. We retain that
cancellation and prove a quantitative comparison with the existing
box and isotropic cutoffs, for the same finite extraction and the
same physical weight. This closes the change-of-cutoff contribution;
it does not assume smallness of the remaining common raw correction.

\subsection{A smooth cutoff for the nonrectangular region}
Choose the function $\chi_0$ in Section~\ref{sec:anisotropic-fixed-count-bands}
also nonincreasing on $[0,\infty)$. Define
\[
 \varphi_0(t)=\chi_0(t),\qquad
 \varphi_j(t)=\chi_0(t/2^j)-\chi_0(t/2^{j-1})\quad(j\ge1).
\]
These functions are nonnegative and sum to one on $\R$. With
$\mathcal D_n,a_n$ from \eqref{eq:adaptive-cross-scales}, put
\begin{equation}\label{eq:cross-smooth-cutoff}
 \Xi_n(v)=\sum_{\mathbf j\in\mathcal D_n}
                 \prod_{i=1}^4\varphi_{j_i}(v_i/a_n),\qquad
 \chi_n^\sharp(\omega)=\Xi_n(\sqrt n\,\omega),\qquad
 K_n^\sharp=\mathcal F^{-1}\chi_n^\sharp.
\end{equation}
For large $n$, the first three physical supports are strictly inside
the chosen torus chart; the expression is extended by zero there.
Its largest coordinate support is $2n^{-41/100}$ and its central
plateau contains the physical ball $|\omega|\le2n^{-99/200}$.
The corresponding rescaled scales are $2n^{9/100}$ and $2n^{1/200}$.

\begin{lemma}[Cutoff plateau and count-coordinate kernel tail]
\label{lem:cross-kernel-tail}
The cutoff satisfies $0\le\Xi_n\le1$, is supported in
$\mathcal H_n$, and is one on the central ball and on the inner
product region \eqref{eq:adaptive-cross-inner}. For every fixed
integer $J\ge0$,
\begin{equation}\label{eq:cross-kernel-Schwartz}
 |K_n^\sharp(x)|\le C_J\sum_{\mathbf j\in\mathcal D_n}
     \prod_{i=1}^4 B_{\mathbf j,i}
          (1+B_{\mathbf j,i}|x_i|)^{-J},\qquad
 B_{\mathbf j,i}=\lambda_i/\sqrt n,
 \quad x\in\Z^3\times\R.
\end{equation}
In particular, for $L_n=\lceil\lambda n\rceil$ with
$\lambda>c_*^{-1}+1$, and the high-count part
$\mathsf T^{a,L_n}$ of the original marked law,
\begin{align}
 n^2\sup_{\mathcal W_{n,R,A}}
        |K_n^\sharp*\mathsf T^{a,L_n}|
 \le C_{A,\lambda,J}M_a n^{41/175-101J/200}.
 \label{eq:cross-count-tail}
\end{align}
The third coordinate, not the largest roof scale, supplies this decay.
\end{lemma}
\begin{proof}
The nonnegative full product partition sums to one, so a subsum lies
between zero and one. Each summand is supported in its box from
\eqref{eq:adaptive-cross-support}. If a summand is nonzero at $v$,
then $\lambda_i\le2\max\{a_n,|v_i|\}$ for every $i$; this is
immediate for $j_i=0$, and follows from
$|v_i|>\lambda_i/2$ for $j_i>0$. At a point of the inner product
region all nonzero summands therefore belong to $\mathcal D_n$,
which proves the plateau. On the central ball only the zero tuple
can contribute and its value is one.

For $j\ge1$, write $\varphi_j(v/a_n)$ as
$\chi_0(v/\lambda_j)-\chi_0(2v/\lambda_j)$. Thus the inverse of
every factor is a dilation of one of two fixed Schwartz functions.
The mixed inverse transform, with the three integer output coordinates
left discrete, gives \eqref{eq:cross-kernel-Schwartz}. Its constants
are uniform over all dyadic indices. The sum is finite at every $n$.

The total variation of the high-count marked measure is at most
$c_*M_a$. On a central window its third-coordinate support is
separated by at least $c_{A,\lambda}n$, by
Lemma~\ref{lem:central-count-support}. Since every
$B_{\mathbf j,3}\ge a_n/\sqrt n=2n^{-99/200}$, each count factor
has decay at most $C n^{-101J/200}$. Multiply the remaining scale
sum by $n^2$ and use \eqref{eq:dyadic-cross-scale-sums}; the result is
$Cn^{41/175}$. This proves \eqref{eq:cross-count-tail}. Arbitrary
inverse powers follow by choosing a fixed $J$ large enough.
\end{proof}

\subsection{The exact signed correction and its cancellation}
Require now the additional finite-record admissibility of the weight
$w_{n,k,R}=c_*a\circ(F_R^*)^k$ from
Definition~\ref{def:finite-record-weight}. Fix one extraction
\[
 \mu^{a,L}=E^{a,L}+Q^{a,L},\qquad
 p^{a,L}=e^{a,L}+r^{a,L},\qquad
 \widehat Q^{a,L}\in L^1(\T^3\times\R).
\]
All cutoffs in this subsection use precisely these same measures.
For a real smooth compactly supported frequency cutoff $\chi$,
with kernel $K_\chi=\mathcal F^{-1}\chi$, define
\begin{equation}\label{eq:coherent-raw-correction}
 \mathcal R_\chi^{a,L}
 =e^{a,L}-K_\chi*E^{a,L}
       +\mathcal F^{-1}[(1-\chi)\widehat Q^{a,L}].
\end{equation}
This is a signed or complex function defined almost everywhere; it
is not the sum of the absolute edge and residual bounds. Its essential
supremum may be infinite. The residual inverse is absolutely convergent,
and the other convolutions are well defined for finite measures.

\begin{proposition}[Exact transport of a raw correction]
\label{prop:exact-raw-cutoff-transport}
For two such cutoffs $\chi,\psi$, almost everywhere on the mixed
space,
\begin{align}
 \mathcal R_\chi^{a,L}&=p^{a,L}-K_\chi*\mu^{a,L},\notag\\
 \mathcal R_\chi^{a,L}-\mathcal R_\psi^{a,L}
             &=(K_\psi-K_\chi)*\mu^{a,L}.
 \label{eq:exact-raw-cutoff-transport}
\end{align}
In particular their difference has a bounded continuous representative,
and
\begin{equation}\label{eq:raw-cutoff-global-comparison}
 \|\mathcal R_\chi^{a,L}-\mathcal R_\psi^{a,L}\|_\infty
 \le(2\pi)^{-4}\int |\chi-\psi|\,|\widehat\mu^{a,L}|\dd\omega.
\end{equation}
The identities hold for any exact finite extraction with the stated
integrable residual transform. They introduce no estimate for $A_2$.
\end{proposition}
\begin{proof}
Absolute residual inversion gives
$\mathcal F^{-1}[(1-\chi)\widehat Q]=r-K_\chi*Q$.
Substitute this in \eqref{eq:coherent-raw-correction} and use
$E+Q=\mu^{a,L}$. Subtract the two identities. In frequency form the
edge change is $(\psi-\chi)\widehat E$ and the residual change is
$(\psi-\chi)\widehat Q$; their sum is
$(\psi-\chi)\widehat\mu^{a,L}$. This is the cancellation that is
lost by taking absolute values of the two changes first. The last
integrand is integrable, since the cutoff difference has compact
support and the measure transform is bounded. Mixed Fourier inversion
therefore proves the norm bound and continuity. The equalities are
almost-everywhere equalities of actual densities, not a subtraction
of two extended-real essential suprema.
\end{proof}

\subsection{Comparison with the already used cutoffs}
Retain $\chi_n^{\rm box},K_n^{\rm box}$ from
\eqref{eq:roof-box-kernel} and
$\chi_n^\natural,K_n^\natural$ from
Section~\ref{sec:rate-preserving-band}. Each is one on the original
central ball for sufficiently large $n$. Define
\begin{equation}\label{eq:raw-cutoff-rate}
 \Delta_n(a)=M_an^{-3/280}+V_an^{-983/350}.
\end{equation}
This exponent for the variation is the slowest of the three
complementary-band estimates, not the weaker central-ball exponent.

\begin{theorem}[A vanishing complete change-of-cutoff contribution]
\label{thm:coherent-raw-cutoff-rate}
For $\chi,\psi$ any two of
$\chi_n^\sharp,\chi_n^{\rm box},\chi_n^\natural$,
$L_n=\lceil\lambda n\rceil$, $\lambda>c_*^{-1}+1$, and fixed
$A,P_0<\infty$ with $P_0>0$,
\begin{align}
 n^2\mathop{\rm ess\,sup}_{\mathcal W_{n,R,A}}
   |\mathcal R_\chi^{a,L_n}-\mathcal R_\psi^{a,L_n}|
       \le C\Delta_n(a)+C_{A,\lambda,P_0}M_an^{-P_0}.
 \label{eq:coherent-raw-cutoff-rate}
\end{align}
The constant is uniform in $R,n,k,a$ in the common weight class.
If in addition $c\lambda>a$ for the constants in
\eqref{eq:count-tail-for-inversion}, the comparison is global, with
an exponentially small tail in place of the last term.
\end{theorem}
\begin{proof}
The cutoff difference vanishes on $|z|\le2n^{-99/200}$ and is
supported in $n^{-1/2}\mathcal U_n^\sharp$. It has modulus at most
one. Equations~\eqref{eq:adaptive-cross-fixed-count},
\eqref{eq:roof-box-fixed-count}, and
\eqref{eq:rate-preserving-annulus} therefore imply
\[
 (2\pi)^{-4}n^2\int |\chi-\psi|\,
                         |\widehat\mu^{[k],a}_{n,R}|\dd z
                       \le C\Delta_n(a).
\]
Here $1/40>3/280$, and the two other variation exponents are larger
than $983/350$. This is an estimate for the full law at the prescribed
return count, not for an averaged coefficient.

In \eqref{eq:exact-raw-cutoff-transport}, write
$\mu^{a,L_n}=\mu^{[k],a}_{n,R}-\mathsf T^{a,L_n}$. The full-law
part has the global bound just proved. The remaining convolution
is bounded on the central window by the two relevant count-separation
bounds: \eqref{eq:cross-count-tail}, \eqref{eq:box-count-separation},
or the isotropic bound in the proof of
Corollary~\ref{cor:rate-preserving-raw-budget}. Choose their fixed
Schwartz orders for $P_0$. This proves the first assertion without
assuming $c\lambda>a$. With that stronger condition the tail total
variation is exponentially small, and all three kernel suprema are
polynomially bounded. Their products are exponentially small globally.
\end{proof}

\begin{corollary}[One common local raw remainder]
\label{cor:one-common-raw-remainder}
Let $\mathcal A_n$ be a class of common admissible weights and marks
for which $\sup_{a\in\mathcal A_n}\Delta_n(a)\to0$ and $M_a$ has
at most fixed polynomial growth. Then
\[
 \sup_{R,k,a\in\mathcal A_n}n^2
        \|\mathcal R_\chi^{a,L_n}\|_{L^\infty(\mathcal W_{n,R,A})}
                   \longrightarrow0
\]
holds for one of the three cutoffs if and only if it holds for all
three. Finite values of the norms likewise transfer, and the difference
of any two finite norms is at most the bound in
\eqref{eq:coherent-raw-cutoff-rate}. For a fixed bounded-BV weight
class the rate is $O(n^{-3/280})$.
\end{corollary}
\begin{proof}
Apply the triangle inequality to the almost-everywhere identity in
Theorem~\ref{thm:coherent-raw-cutoff-rate} in both directions, choosing
$P_0$ larger than the polynomial growth exponent of $M_a$. If one norm
is finite, the bounded difference makes the other finite. No difference
of infinite norms is formed. The uniform convergence statement follows.
\end{proof}

The corollary does not transfer smallness of the isolated edge term or
of the absolute complementary integral. Cancellation between them is
retained in $\mathcal R$. In particular a small value of their signed
sum cannot be used to bound either of the two absolute quantities.
This distinction permits one fixed raw problem to be used throughout
the expansion of the known Fourier region.

\subsection{The exact raw budget at the new cutoff}
For the new kernel define, immediately before use,
\begin{align}
 \mathcal E^{\sharp,a,L}_{n,k,R,A}
 &=\mathop{\rm ess\,sup}_{\mathcal W_{n,R,A}}
                         |e^{a,L}-K_n^\sharp*E^{a,L}|,\notag\\
 \mathcal C^{\sharp,a,L}_{n,k,R}(B)
 &=\int_{\T^3\times[-B,B]}
           |1-\chi_n^\sharp|\,|\widehat Q^{a,L}|\dd\omega.
 \label{eq:cross-raw-terms}
\end{align}
The first quantity is allowed to be infinite. The exact identity on
the central count labels is
\begin{equation}\label{eq:cross-raw-exact-identity}
 p_{n,R}^{[k],a}=K_n^\sharp*\mu_{n,R}^{[k],a}
                  +\mathcal R_{\chi_n^\sharp}^{a,L_n}
                  -K_n^\sharp*\mathsf T^{a,L_n}.
\end{equation}
It follows from central count support and the first identity in
Proposition~\ref{prop:exact-raw-cutoff-transport}.

\begin{corollary}[Raw inversion with the coherent cutoff]
\label{cor:cross-raw-budget}
For $B\ge1$, $P_0>0$, and
$y=((\ell,t)-n\bar G_R)/\sqrt n$,
\begin{align}
 &\mathop{\rm ess\,sup}_{\mathcal W_{n,R,A}}
 |n^2p_{n,R}^{[k],a}(\ell,t)-\alpha_a\mathfrak g_R(y)|\notag\\
 &\quad\le C\left(M_an^{-3/280}\sqrt{\log(2+n)}
                              +V_an^{-9/175}\right)
        +CM_ae^{-cn^{1/100}}+C_{A,\lambda,P_0}M_an^{-P_0}\notag\\
 &\qquad+n^2\mathcal E^{\sharp,a,L_n}_{n,k,R,A}
       +(2\pi)^{-4}n^2\mathcal C^{\sharp,a,L_n}_{n,k,R}(B)
       +\frac{n^2}{\pi B}A_2(n,L_n,R,w_{n,k,R}).
 \label{eq:cross-raw-budget}
\end{align}
The inequality has its extended-real meaning if necessary.
\end{corollary}
\begin{proof}
Apply \eqref{eq:adaptive-cross-central} to the first term of
\eqref{eq:cross-raw-exact-identity}, using $0\le\Xi_n\le1$.
The omitted Gaussian frequencies are outside the plateau containing
$|v|\le2n^{1/200}$; uniform ellipticity gives the exponential term.
Use \eqref{eq:cross-count-tail} for the high-count convolution.
Bound the two parts of the coherent correction separately only now:
its edge term is the one in \eqref{eq:cross-raw-terms}. Splitting the
absolutely convergent residual inverse at $|b|=B$ and using
\eqref{eq:finite-count-far-tail} gives the last two terms. This keeps
the true derivative sum at $L_n\asymp n$ unchanged.
\end{proof}

\paragraph{The remaining absolute and conditional estimates.}
Theorem~\ref{thm:coherent-raw-cutoff-rate} is an unconditional estimate
for the entire signed change of the raw correction, and
Corollary~\ref{cor:one-common-raw-remainder} makes that comparison
independent of which of the three cutoffs is used. It supplies neither
an absolute value for the common correction nor a bound on the genuine
long-time $A_2$. Frequencies outside the retained union, the common
local raw remainder, weighted raw denominators and relative physical
event replacement still require estimates for the same original law.
The last comparison concerns the same indicator on both sides. No
physical observation event has been substituted.
